\section{Cierre y Fuentes}

% ================================================================
\begin{frame}{Ideas Clave que Deben Quedar}
\footnotesize
\begin{columns}[T,onlytextwidth]
\begin{column}{0.62\textwidth}
\begin{enumerate}\setlength{\itemsep}{0.22em}
  \item \textbf{Aprendizaje no paramétrico}: KNN no asume una función global $\hat{f}(X)$; aproxima la densidad condicional evaluando vecindades locales $\mathcal{N}_k(x_0)$.
  \item \textbf{La métrica define la vecindad}: cambiar la distancia (Euclidiana, Manhattan, Chebyshev) altera la geometría de las regiones de decisión.
  \item \textbf{Escalamiento obligatorio}: diferencias de escala dominan la distancia; usar siempre un \texttt{Pipeline} con \texttt{StandardScaler}.
  \item \textbf{Hiperparámetro $k$}: controla el trade-off sesgo--varianza ($k=1$ sobreajusta; $k$ grande suaviza en exceso).
  \item \textbf{Maldición de la dimensionalidad}: en alta dimensión los puntos se vuelven equidistantes y el soporte local deja de ser representativo.
  \item \textbf{Lazy learner}: entrenamiento trivial $O(1)$, pero costo computacional desplazado a la inferencia $O(N \cdot p)$.
\end{enumerate}
\end{column}

\begin{column}{0.35\textwidth}
\begin{keybox}
\centering
\textbf{Vecindad}\\[0.08cm]
$+$\\[0.08cm]
\textbf{Métrica}\\[0.08cm]
$+$\\[0.08cm]
\textbf{Escala}\\[0.10cm]
$\Downarrow$\\[0.10cm]
\scriptsize
\textbf{Fronteras flexibles sin supuestos paramétricos}
\end{keybox}

\vspace{0.10cm}
\centering
\tagpill{Distancia} \quad \tagpill{Escala} \\[0.10cm]
\tagpill{Sesgo-Varianza} \quad \tagpill{Pipeline}
\end{column}
\end{columns}
\source{James et al. (2013); Hastie et al. (2009); Géron (2022).}
\end{frame}

% ================================================================
\begin{frame}{Bibliografía Base Recomendada}
\footnotesize
\begin{ccdebox}
\textbf{James, G.; Witten, D.; Hastie, T.; Tibshirani, R.} (2013). \emph{An Introduction to Statistical Learning with Applications in R}. Springer.\\
\textbf{Capítulos 2 y 4:} fundamentos de aprendizaje supervisado, clasificador de Bayes, formulación de KNN y trade-off sesgo--varianza.
\end{ccdebox}
\vspace{0.06cm}
\begin{ccdebox}
\textbf{Hastie, T.; Tibshirani, R.; Friedman, J.} (2009). \emph{The Elements of Statistical Learning}, 2.ª ed., Springer.\\
\textbf{Capítulos 2 y 13:} teoría de estimadores locales por densidad de kernel, prototipos, tasas de error y maldición de la dimensionalidad.
\end{ccdebox}
\vspace{0.06cm}
\begin{ccdebox}
\textbf{Géron, A.} (2022). \emph{Hands-On Machine Learning with Scikit-Learn, Keras, and TensorFlow}, 3.ª ed., O'Reilly.\\
\textbf{Capítulo 3:} implementación práctica de \texttt{KNeighborsClassifier}, escalamiento y validación con \texttt{Pipeline} y \texttt{GridSearchCV}.
\end{ccdebox}
\end{frame}
